\section{The isotropic improvement}\label{sec:isotropic}

This section treats the one-spatial-coordinate model and the isotropic
two-spatial-coordinate model. Its purpose is to turn incidence
information along trajectories into a quadratic test with a smaller
thickness and the same time interval. We first construct a velocity
graph, prove that its coefficients have polynomial fits along many
trajectories, and interpolate these fits to a polynomial field.
A switched-line argument makes that field approximately conservative.
The final step restores the degree allowed for a chart. In two
spatial coordinates, broad directions can still lie on a conic; the
higher-order terms they fail to detect are expressed through the
hidden coordinate to obtain a quadratic test.

\begin{theorem}[Isotropic improvement]\label{thm:isotropic-improvement}
Consider the tangent data of \Cref{prop:tangent-palette}, with
\(0<k\leq1\) and \(0\leq d<1\), either in version~1, where \(j=1\), or
in version~2 with \(j=2\) and optimized narrowness \(\ell=0\).
In every substantial residual of the reference path one can, after
passing to a subsequence, find a packet \(Q=Q_s\), of thickness
\(a=N^{-s}\) and time length \(h\sim_{\log,N}a^k\), and a
time-improvement candidate of thickness
\[
 a_2=ah^{e_2},\qquad e_2>0.
\]
The candidate has incidence mass at least
\[
 \frac{h\nu_Q}{K}\left(\frac{a_2}{a}\right)^d,
 \qquad K=N^{o(1)}.
\]
Writing the candidate as \(P_{\rm new}\), its bounds
\(|P_{\rm new}|\leq Ka_2\) and
\(|\partial_wP_{\rm new}|\leq K\) hold on its assigned
trajectories throughout the same \(h\)-length block. The hidden
derivative is bounded below by \(K^{-1}\) on the counted incidences,
and its hidden curvature is at most \(K/a_2\). In version~1 the test
has the special form \(w-F(t,y)\), with \(F\) quadratic. In version~2
it is a quadratic in \((t,y,w)\).
\end{theorem}

Here \emph{substantial} refers to positive probability in the
original reference path, whereas a dense discovery family may have
only inverse-subpower probability. The distinction is needed when
\Cref{prop:candidate-upgrade} converts candidates into a known
atlas. We do not use an arbitrary tangent subpower as an actual
\(N_0\)-subpower bound for an output.

\subsection{Measures, scales, and inherited estimates}

Write
\[
 z=(t,y),\qquad m=1+j,\qquad {\bf v}=(1,V),\qquad W=(z,X).
\]
A parent world and its label are included in every state below.
The parent prior is \(\nu\), and restrictions of the incidence measure
are denoted by \(\lambda\). They are generally left unnormalized:
\begin{equation}\label{a05:incidence-density}
 d\lambda(l,t)=G_{\rm wt}(l,t)\,d\nu(l)\,dt,
 \qquad 0\leq G_{\rm wt}\leq K.
\end{equation}
All initial losses are summed in the path mixture of parents.
If a line portion has weighted duration at least \(H/K_1\),
its preceding incidence mass bounds its prior weight times \(H/K_1\);
its ordinary duration is at least \(H/(KK_1)\).
We use this before further restrictions, rather than assume that a
later restricted family still saturates a prior mass.

All exponent ranges are fixed first, with room for finer
comparisons. A finite list of pure and joint caps is prepared
before sparse selection. Lists of exponents subsequently increase
slowly enough that products of their subpower losses remain
subpower. Deletion budgets are assigned for the whole finite list
before it is enlarged. We use successively slower grids for
physical counts, parameter estimates, field comparisons, strip
filling, and residual comparisons. Interpolation between neighboring
dyadic widths whose ratio is at most \(K\) is performed only after
the endpoint estimates. Its losses are not repeatedly multiplied
within the grid they complete.

We use the following consequences of
\Cref{prop:tangent-palette,a03:finite-state-conditioning}.
For every required fixed power width \(p\), the scalar cap in \(X\)
is \(Kp^d\) per unit \(z\)-volume. The velocity palette \(\pi\)
has a positive angular exponent. Joint upper bounds have the same
pure cap multiplied by \(K\pi(J)\), for a velocity bin \(J\).
They remain valid for later submeasures relative to the earlier
state weights. Conditional palette domination on a later dense
family follows after deleting states with insufficient denominators.
There is no assertion of velocity independence at several times.

The true packets \(Q_u\), at thickness \(p=N^{-u}\) and time length
\(H\sim_{\log,N}p^k\), have
\begin{equation}\label{a05:packet-mass}
 \nu_Q\sim_{\log,N}p^dH^j.
\end{equation}
Their assigned trajectory sets are disjoint in each time block,
and assignments are invariant throughout that block. All spatial
axes have length \(H\) up to \(K\), and center velocities are
bounded by \(K\). In normalized packet coordinates the base
incidence density is at most \(KH\nu_Q\).
The stable graph lists in the tangent preparation have at most
\(K\) ball/sign entries per packet. Their branches are holomorphic
algebraic functions of bounded relation degree on complex
enlargements of their real comparison balls.
We use \Cref{lem:algebraic-norm} only on interiors with these margins.

\subsection{Physical support and parameter masses}

The physical support will have exponent \(m+d\), strictly below
its ambient dimension \(m+1\). This deficit will force the lifted
velocities \(({\bf v},X')\) near an \(m\)-dimensional graph.
We also need bounds for the masses of trajectory coefficient bins.
Their lower bounds provide the denominators needed to retain
controlled spatial caps after conditioning and rescaling.

We first prove, after negligible deletion on the required grids,
\begin{equation}\label{a05:physical-counts}
 \#W_p\leq Kp^{-(m+d)},\qquad
 \lambda(W_p)\leq Kp^{m+d},\qquad
 \lambda(z_p)\leq Kp^m.
\end{equation}
Here \(\#W_p\) counts occupied cubes in a parent. The two mass
ceilings follow from the pure caps. The support argument requires
a gradient bound for the packet graphs.

At scale \(p\), each packet has at most \(K\) local predictions
\[
 X=f(z)+O(Kp).
\]
In version~2 these are obtained by composing the native quadratic
test with a simple root of the packet equation, or with its affine
quadratic center in the large-curvature case. The root error is
\(O(Kp)\) in the hidden coordinate; the native derivative and
curvature bounds make the error in \(X\) also \(O(Kp)\).
Comparison balls have inverse-subpower radii in coordinates
scaled by \(H\), and can be assigned from the exact base by a
lattice, with at most two root signs. Removing entries of incidence
mass below \(H\nu_Q/K_1\) costs at most \(K/K_1\), since
\(\sum_QH\nu_Q\leq1\). The remaining real projections have relative
volume at least \(K^{-1}\) by the packet base ceiling.
Thus \Cref{lem:algebraic-norm} bounds all fixed derivatives of
these predictions by \(K\) in the scaled coordinates.

Prune line/packet/ball/sign portions of weighted duration below
\(H/K_1\), at total cost at most \(K/K_1\).
On a surviving line, \(f(z_l(t))-X_l(t)\) is holomorphic algebraic
of bounded relation degree, and bounded by \(Kp\) on a time set
of length at least \(H/K\). Its chord lies in the complex
comparison enlargement with a fixed relative margin.
The one-variable form of \Cref{lem:algebraic-norm} gives
\begin{equation}\label{a05:graph-directional}
 |D_{\bf v}f-X'|\leq Kp/H.
\end{equation}
As \(k\leq1\) and \(|X'|\leq K\), the directional derivative is
bounded by \(K\).

If \(|\nabla_zf|\) were power-large on a substantial subfamily,
\Cref{a05:graph-directional} and \(|V|\leq K\) would make its
spatial part power-large, and
\[
 \partial_tf+\nabla_yf\cdot V=O(K)
\]
would predict \(V\) in a power-thin affine strip (an interval for
\(j=1\)). The state consisting of packet, ball, sign and a fine
base bin has actual list knowledge at \(p_L\) and only \(K\)
ambiguity from a deep end history. The Hessian costs at most
\(K/H^2\), so the base bin can freeze the gradient to error one.
This contradicts \Cref{a03:breadth} in version~2, or the
angular conclusion of \Cref{prop:tangent-palette} in version~1.
Sending fixed-power tolerances slowly to zero allows
\(|\nabla_zf|\leq K\) outside negligible mass. Remove light graph
entries again. Their remaining base volume and
\Cref{lem:algebraic-norm}, applied to physical derivative
components, extend this bound throughout the used interiors.

A packet now uses at most \(K(1+H/p)^m\) physical \(W_p\)-cubes.
Since \(H\geq p/K\), summing with \Cref{a05:packet-mass} over
packets and \(H^{-1}\) blocks gives
\[
 \frac{K}{H}\,\frac{1}{p^dH^j}
       \left(\frac Hp\right)^m
 \leq Kp^{-(m+d)}.
\]
This proves \Cref{a05:physical-counts}. Remove globally light
\(W_p\)-cubes of mass below \(p^{m+d}/K_1\); a sufficiently large
subpower \(K_1\) makes the loss negligible. The old ceilings then
give the hereditary support bounds
\begin{equation}\label{a05:local-support}
 \#\{W_p\subset W_{w_0}\}\leq K(w_0/p)^{m+d}\quad(w_0\geq p),
 \qquad
 \#\{W_p\text{ over a fixed }z_p\}\leq Kp^{-d},
\end{equation}
with bounded enlargements at boundaries.

Let \(B_p\) be an aligned dyadic cube in all coefficients of
\[
 X=x_0+tx_1+t^2x_2,\qquad y=y_0+tV.
\]
Additional trajectory tags stay in the prior but are not part of
this coefficient bin. There is one fixed, unnormalized,
preceding long-line prior \(\nu^\circ\) for which every
participating block satisfies
\begin{equation}\label{a05:parameter-floor}
 K^{-1}p^{j+d}\pi([V]_p)
 \leq \nu^\circ(B_p)
 \leq Kp^{j+d}\pi([V]_p).
\end{equation}
A physical cell \(W_p\) and velocity bin \([V]_p\) admit at most
\(K\) participating \(B_p\)'s. We now prove these assertions,
keeping prior and incidence denominators separate.

The joint caps imply, for a time interval \(I\), a \(y\)-box
of volume \(J_y\), an \(X\)-interval of width \(w_0\), and a
velocity bin \(J\),
\begin{equation}\label{a05:rectangular-cap}
 \lambda(\text{these constraints})
       \leq K|I|J_yw_0^d\pi(J).
\end{equation}
For the current finite grid of spatial meshes, take one common time
partition much finer than every mesh on that grid.
Deleting line/bin pairs of weighted duration below \(|I|/K_1\)
costs at most \(1/K_1\). Bounded speed makes a constraint at one
time hold with a mesh enlargement throughout \(I\).
Here a line is \emph{active at \(t\in I\)} when its whole line--\(I\)
mark survives this duration test. Activity does not condition the
trajectory prior on incidence success at the particular time \(t\).
Dividing \Cref{a05:rectangular-cap} by the retained duration, as in
\Cref{a03:retained-duration-cap}, gives the fixed-time active-prior estimate
\begin{equation}\label{a05:fixed-time-cap}
 \nu\{l:\text{active at }t,\ (y_l(t),X_l(t))\in B,\
                      [V_l]_\beta=J\}
       \leq Kp^{j+d}\pi(J)
\end{equation}
for a spatial cube \(B\) of side \(p\); the rectangular version
holds too. Remove null exceptional times and lines with tiny
total weighted occupation. Restrict \(\nu\), without normalization,
to lines of occupation at least \(1/K\); this defines \(\nu^\circ\).
Freeze this prior after the common finite-grid preparation.
Later labels satisfy \(\lambda\leq K\nu^\circ\,dt\); restricting
them never renormalizes \(\nu^\circ\).

For an affine bin \(\alpha=[y_0,V]_p\), let
\(\Pi_\alpha=\pi([V]_p)\). Integrating on the preceding long
occupation gives
\[
 \nu^\circ(\alpha)\leq Kp^j\Pi_\alpha,\qquad
 \nu^\circ\{\alpha,\ (x_0,x_1,x_2)\in B_r\}
       \leq Kr^dp^j\Pi_\alpha\quad(r\geq p).
\]
At each time these coefficient restrictions put the position in
the corresponding moving boxes, so
\Cref{a05:rectangular-cap} applies. This proves the upper half of
\Cref{a05:parameter-floor}.

If the lower half failed by a fixed power on a substantial
residual, a parent and an \(\alpha\) would have bad incidence mass
at least \(p^j\Pi_\alpha/K\). This follows by summation, because
the sum of \(p^j\Pi_\alpha\) over affine bins is at most \(K\);
ignore null palette bins. Restrict \(\nu^\circ\) to the bad
coefficient blocks in this \(\alpha\), and divide it by
\(p^j\Pi_\alpha\). The scalar coefficient-ball caps are \(Kr^d\),
the bad labels have dense mass, and their \((t,X)\)-support has at
most \(Kp^{-1-d}\) cells: \(\alpha\) fixes \(y\) to \(K\) options
per time cell, and \Cref{a05:local-support} counts scalar options.
\Cref{lem:scalar-clustering} gives a scalar coefficient \(p\)-bin
of mass at least \(p^d/K\). Splitting among boundedly many aligned
bins if needed contradicts the fixed-power deficiency of every bad
\(B_p\). Slowly relaxing the deficiencies proves the lower bound.
For \(p^{-1}\leq K\), delete light bins directly by their subpower
count.

A block \(B_p\) uses at most \(K/p\) physical cells \(W_p\).
Discard entries below \(p\nu^\circ(B_p)/K_1\), at total cost
\(K/K_1\). For a fixed cell and velocity bin their summed mass is
at most \(Kp^{m+d}\pi([V]_p)\), by the joint cap, and every
surviving entry has mass at least \(p^{m+d}\pi([V]_p)/K\).
There are at most \(K\) such blocks. In particular,
\begin{equation}\label{a05:derivative-list}
 X'\text{ to width }p\text{ has at most }K\text{ options given }
                    (W_p,[V]_p).
\end{equation}

Between grid widths \(p'<p''\) with \(p''/p'\leq K\), discard
child velocity bins whose palette mass is too small relative to
the containing bin. Each parent has at most \(K\) children, so
palette domination makes the loss negligible. The finer lower
bound implies intermediate lower bounds by inclusion; intermediate
upper lists use finer lists and at most \(K\) refinements.
For upper prior and fixed-time estimates, sum the finer palette
factors. All of \Cref{a05:parameter-floor} concerns the same
\(\nu^\circ\); later restrictions do not assert new incidence
saturation.


\subsection{A velocity graph and its coarse lists}

Fix the scales of the intended improvement before constructing its
graph. Choose \(a=N^{-s}\), \(s>0\), its actual packet time
\(h\sim_{\log,N}a^k\), and an integer \(M\geq2\) with
\(h^M\ll a^2\) by a fixed power. Choose the scalar-normal floor
\(\zeta\) sufficiently fine for \Cref{prop:higher-jets} at \(h,M\),
and then choose
\[
 \rho\ll\zeta,\qquad \rho\ll h^{M+1/4}.
\]
The exponent \(\kappa>0\) in the next proposition depends only on
\(m,d\). Take \(b\) so fine that \(b/\rho\) and
\(b^\kappa/\rho\) are negligible at every required comparison scale.
The finite interpolation meshes and descendant graph scales are
included in this choice, with fixed power margins. These choices
depend on the fixed exponents and \(M\), before any test of failure
for the graph; none will be changed inside such a failure set.

On any substantial residual, renew breadth at the fine physical states
\(W_b\). Outside negligible mass, each conditional velocity law puts
as small a fixed probability as desired in any fixed-power thin affine
strip, or interval if \(j=1\). Letting that power tend slowly to zero
gives reciprocal-subpower separation. Thus \(m\) independent draws
from a state span with determinant at least \(K^{-1}\), with
probability bounded below. For \(j=2\), choose separated points and
then avoid their joining lines; for \(j=1\), choose a separated pair.
The same assertion holds for unit directions. The probability remains
bounded below after sufficiently small fixed fractional losses in
each state. Conditional domination by \(K\pi\) is prepared at the same
time by deleting light states, retaining the unnormalized bounds
against the earlier state weights.

We use the determinant multilinear Kakeya estimate in the following
form. If \(\mathcal T_i\), \(1\leq i\leq n\), are weighted families of
radius-\(r\) neighborhoods of full affine lines in \(\R^n\), with
unit directions \(e_T\), nonnegative weights \(w_T\), and
\(M_i=\sum_{T\in\mathcal T_i}w_T\), then
\begin{equation}\label{a05:determinant-kakeya}
 \int_{\R^n}
 \left(
 \sum_{\substack{T_i\in\mathcal T_i\\1\leq i\leq n}}
 |\det(e_{T_1},\ldots,e_{T_n})|
 \prod_{i=1}^n w_{T_i}\1_{T_i}(x)
 \right)^{1/(n-1)}dx
 \leq C_n r^n\prod_{i=1}^n M_i^{1/(n-1)}.
\end{equation}
This is the equal ambient and wedge dimension case of
\citet[Theorem~1]{CarberyValdimarsson}, after scaling unit-radius
tubes by \(r\).
The near-endpoint multilinear theorem is due to
\citet{BennettCarberyTao2006}; \citet{Guth2010} proved the endpoint,
and \citet[Section~7]{BourgainGuth2011} obtained a wedge-weighted
formulation.
Only \(2\leq n\leq4\) is used. The statement for line measures follows
by rounding line parameters arbitrarily finely on a bounded region,
summing their weights, enlarging tubes by a bounded factor, and using
slightly stronger determinant thresholds before passage to the limit.
No cardinality bound for the rounded lines is required.

\begin{proposition}[Physical velocity graph]\label{a05:velocity-graph}
For a sufficiently fine fixed power \(b\), there is a row
\(G=(G_0,G_y)\), constant on each \(W_b\) and bounded by \(K\), such
that on a substantial reference family
\begin{equation}\label{a05:velocity-graph-estimate}
 |X'(t)-G(W_b)\cdot{\bf v}|\leq b^\kappa,
\end{equation}
where \(\kappa>0\) depends only on \(m\) and \(1-d\).
For every required dyadic \(s_0\geq b^\kappa\), the values of \(G\)
to width \(s_0\) have at most \(K\) possibilities in a fixed
coarse physical cell \(W_{s_0}\).
\end{proposition}

\begin{proof}
Put \(n'=m+1\), choose
\[
 0<\theta<(1-d)/3,\qquad 0<\tau<\theta/(n'-1).
\]
If no \(m\)-plane contains, say, \(0.9\) of the unit lifted directions
\(({\bf v},X')\) within distance \(b^\tau\), successive independent
samples have distances at least \(b^\tau\) from the preceding spans
with probability bounded below. Enlarge a smaller span to an
\(m\)-plane when applying the hypothesis. The resulting
\(n'\)-fold determinant is at least a constant times
\(b^{\tau(n'-1)}\), and hence is at least \(b^\theta\) for large \(N\).

Suppose these bad cells carry substantial mass. Group them into
dyadic \(W\)-cubes of side \(w_0\sim b^{2/3}\). Delete line portions
with weighted duration below \(w_0/K_1\) in such a cube. A bounded
quadratic trajectory meets at most \(K/w_0\) cubes, so this costs
\(K/K_1\). Remove bad cells with excessive fractional loss and
cubes retaining too little of their preceding bad mass.
In a remaining cube of preceding mass \(M>0\), the retained cell
masses \(m_E\) sum to at least a fixed fraction of \(M\), determinant
success still has probability bounded below, and the contributing
prior weight is at most \(KM/w_0\).

Approximate each trajectory there by its affine tangent at the cube's
central time. The position error is \(O(Kw_0^2)\ll b\), and the
direction error is \(O(Kw_0)\ll b^\theta\).
Radius-\(O(b)\) tubes about these full tangent lines contain every
successful cell and retain its determinant threshold.
A trajectory contributes at most \(Kb\) weighted duration to a
cell. Therefore the successful tuples have prior product weight
at least \((m_E/b)^{n'}/K\).
Apply \Cref{a05:determinant-kakeya} with \(n=n'\) and cancel the
cell-volume factor \(b^{n'}\):
\[
 \sum_E\left(
 b^\theta(m_E/b)^{n'}/K
 \right)^{1/(n'-1)}
 \leq K(M/w_0)^{n'/(n'-1)}.
\]
The power-mean inequality forces at least
\(K^{-1}b^\theta(w_0/b)^{n'}\) cells.
By \Cref{a05:local-support} there are at most
\(K(w_0/b)^{m+d}\). The ratio of the asserted lower and upper
bounds is at least
\[
 K^{-1}b^{\,\theta-(1-d)/3}\longrightarrow\infty,
\]
a contradiction.

Thus the lifted directions concentrate near an \(m\)-plane on
almost all the relevant cells. Even within the captured fraction,
breadth supplies \(m\) base directions with determinant at least
\(K^{-1}\). Inverting this basis and using bounded \(X'\) shows that
a unit normal to the plane has final component at least \(K^{-1}\).
The plane is therefore a graph with row bounded by \(K\).
Restricting to its captured incidences proves
\Cref{a05:velocity-graph-estimate}, for instance with \(\kappa=\tau/2\).

Keep \(G\) fixed and renew spanning and palette domination at
\(W_b\) on the captured family. Name the resulting incidence
submeasure in parent \(\gamma\) by \(\lambda_{{\rm ref},\gamma}\),
and retain the prepared parent-mixture weights \(p_\gamma\). Thus
\begin{equation}\label{a05:graph-reference}
 \Lambda_{\rm ref}=\sum_\gamma p_\gamma
             \lambda_{{\rm ref},\gamma},\qquad
 L_{\rm ref}=\Lambda_{\rm ref}(\mathrm{all}).
\end{equation}
This reference includes trajectory indices and true packet labels.
It is substantial: the bad-cell mass tends to zero, good cells
retain their fixed captured fraction, and the renewal has vanishing
loss. Hence \(L_{\rm ref}\) is bounded below by a positive constant
on the working subsequence. Freeze this post-renewal reference
together with \(G\).

For each fine cell, the tuples of width-\(s_0\)
velocity bins supporting actual spanning witnesses have product
palette mass at least \(K^{-1}\). For a fixed tuple,
\Cref{a05:derivative-list} and
\Cref{a05:velocity-graph-estimate} allow at most \(K\) row bins.
When \(s_0\) is sufficiently small relative to the reciprocal
conditioning bounds, invert at the bin centers; otherwise the
bound follows directly from \(|G|\leq K\).
Choose one fine-cell witness for each row bin occurring in a
coarse \(W_{s_0}\). Sum its product-palette mass: every tuple pays
for at most \(K\) row bins and the total product palette is one.
There are at most \(K\) row bins. This is a support statement on
the reference family, and survives every later restriction.
\end{proof}

\subsection{Polynomial fits along reference trajectories}

Use the scales \(a,h,M,\zeta,\rho,b\) fixed before the graph
construction. The next assertion concerns its frozen incidence
law \(\Lambda_{\rm ref}\), rather than a later discovery family.

\begin{proposition}[Polynomial qualification]\label{a05:qualification}
There are a subpower \(K_{\rm qual}=N^{o(1)}\) and sets \(U_N\) of
graph-reference incidences such that
\begin{equation}\label{a05:qualification-exception}
 e_N:=\frac{\Lambda_{\rm ref}(U_N)}{L_{\rm ref}}
       \longrightarrow0.
\end{equation}
Every reference incidence outside \(U_N\) admits a row-valued
polynomial in \(t\), of degree at most \(M\) and coefficients
bounded by \(K_{\rm qual}\), which approximates \(G\) within
\(h^M\) on reference-labeled times of its indexed trajectory of
ordinary duration at least \(h/K_{\rm qual}\) in the same
\(Q_s\)-block, including that incidence.
\end{proposition}


The coarse lists do not yet compare the rows at nearby points: two
points in the same cell may use different members of its list. The
following selection makes a common choice in every retained cell. We
will use it first for the physical row and then for a scalar residual
inside a parameter block.

\begin{lemma}[A common choice from local lists]
\label{a05:common-local-choice}
Let \(\mu\) be an unnormalized finite measure on a finite collection of
worlds \(\omega\). In each world let \(Z(e)\in\mathbb R^D\) and
\(Y(e)\in\mathbb R^q\) be measurable functions of an incidence \(e\),
with \(|Y(e)|\leq B\). Fix an integer \(L\geq2\), bounds
\(K_1,\ldots,K_L\geq1\), and tested widths
\[
 1=s_1>s_2>\cdots>s_L=r_*,\qquad
 A=\max_{i<L}s_i/s_{i+1}.
\]
At width \(s_i\), suppose that in every half-open \(s_i\)-cell of
\(Z\) in each world, the occurring values of \(Y\) meet at most \(K_i\)
half-open \(s_i\)-cells. There is a restriction \(\mu'\leq\mu\) such
that, in every world,
\[
 \mu'(\omega)\geq
 \mu(\omega)\prod_{i=1}^L(4^DK_i)^{-1},
\]
and any two retained incidences in the same world satisfy
\[
 |Y(e)-Y(\widehat e)|
 \leq C_{D,q}(A+B)
       \bigl(|Z(e)-Z(\widehat e)|+r_*\bigr).
\]
If, within each world, both \(Z\) and \(Y\) are determined by a
discrete state, the restriction keeps or removes each whole state
of that world. No measure is renormalized.
\end{lemma}

\begin{proof}
Order the scales as displayed. At scale \(s_i\), color the coordinate
cell \(s_i(n+[0,1)^D)\) by \(n\bmod4\). In each world retain a color
of largest current mass. In each retained cell choose a \(Y\)-cell
of largest current mass among its at most \(K_i\) possibilities,
and retain its incidences. Fix a lexicographic rule for ties.
The color depends only on the world and the scale; the value-cell
choice depends only on the world, the scale, and the coordinate
cell. In particular two incidences in the same retained coordinate
cell use the same value-cell choice.

These two operations keep at least \((4^DK_i)^{-1}\) of the current
mass in each world. Their successive application proves the mass
bound. All lists may be frozen before selection, since restrictions
preserve upper lists even when they remove the original witnesses.

Distinct coordinate cells of the same color are separated by at
least \(3s_i\) in one coordinate. Thus
\(|Z(e)-Z(\widehat e)|_\infty\leq s_i\) forces the same cell,
and the selected value-cell gives
\(|Y(e)-Y(\widehat e)|\leq\sqrt q\,s_i\).
For coordinate distance at most one, take the tested width just
above the larger of that distance and \(r_*\). This width is at
most \(A\) times that larger quantity. For larger distance use
\(|Y(e)-Y(\widehat e)|\leq2B\). Norm equivalence proves the stated
comparison. When the two maps are state functions, every test is
a test of the whole state.
\end{proof}

We apply this lemma only on fixed finite grids first. When all
entering masses are inverse-subpower and \(B,K_i=N^{o(1)}\), choose
successively larger grids slowly enough that
\[
 DL\log4+\sum_i\log K_i=o(\log N),\qquad A=N^{o(1)}.
\]
The selected mass is then inverse-subpower and the comparison costs
only \(K\). To obtain this diagonal, at each fixed finite-grid stage
include the entering density and all preceding preparation costs.
Delay stage \(r\) until their combined logarithmic cost is at most
\((\log N)/r\). Let the maximal gap between the tested exponents tend
to zero on the same diagonal. There is one selection loss per
scale, irrespective of the number of cells or worlds. Intermediate
widths require only the comparison above, with no further selections.

\begin{proof}[Proof of \Cref{a05:qualification}]
Fix \(\epsilon>0\). Let \(U_{N,\epsilon}\) be the reference
incidences for which no such polynomial has coefficient bound
\(N^\epsilon\), error at most \(h^M\), and witness duration at
least \(hN^{-\epsilon}\). Use finite coefficient grids with total
evaluation error at most \(h^M/4\) on the bounded parent-time
interval. The discovery below has an extra factor \(h^{1/4}\),
leaving room for this rounding and for the coefficient bound.

If \(\Lambda_{\rm ref}(U_{N,\epsilon})/L_{\rm ref}\) does not
tend to zero, pass to a subsequence where it is bounded below.
Since \(L_{\rm ref}\) is bounded below, this is a substantial
residual of the original reference path with its unchanged parent
weights. We shall find qualifying occurrences inside this residual.
Keep \(G\) and all its mesh choices fixed, and renew breadth on
\(W_b\) in the residual. Until the
strip-filling step below, selections keep whole \(W_b\)'s or remove
only sufficiently small fixed fractions of their laws.

We will apply \Cref{prop:higher-jets} inside a full coefficient
block \(B_\rho\). After subtracting its central curves and dividing
positions by \(\rho\), a comparison of nearby values of \(G_y\)
will supply a time-only row. To obtain scalar-normal data for the
fixed prior \(\nu^\circ(\,\cdot\mid B_\rho)\), we must also
control the scalar support, the spatial masses conditional on smaller
coefficient bins, and the residual velocity. The construction below
establishes these bounds before restricting to one such world.

Apply \Cref{a05:common-local-choice} on the grid between \(\rho\)
and one, with parent as world, \(Z\) the center of \(W_b\), and
\(Y=G(W_b)\). The coarse lists in \Cref{a05:velocity-graph} provide
the hypotheses; cells meeting a coarse-cell boundary require only
bounded enlargements. This retains dense mass and keeps whole fine
states, so their conditional velocity distributions are unchanged.
The displacement from an incidence to its state center is \(O(Kb)\),
which is negligible compared with \(\rho\). Hence the retained
incidences obey
\begin{equation}\label{a05:field-lipschitz}
 |G(W_b)-G(\widehat W_b)|
       \leq K(|W-\widehat W|+\rho).
\end{equation}

In a physical cube \(C=W_\rho\), choose a row \(G_C\) within
\(K\rho\) of all its active rows. For each required
\(1\geq s'\geq\zeta\), set \(p=\rho s'\) and assign whole
\(W_b\)'s by their centers to width-\(p\) strips in
\(X-G_C\cdot z\).
At active times the normal derivative is \(O(K\rho)\).
It is affine along a trajectory and remains so bounded between
any two of those times. The normal range inside \(C\) is at most
\(K\rho^2+Kb\ll p\), so a trajectory meets at most \(K\) strips
there, and at most \(K/\rho\) cube/strip pairs in total.

Delete line/strip portions of duration below \(\rho/K_1\);
remove whole \(W_b\)'s with excessive fractional loss, then
strips with excessive total loss. Each surviving strip of
preceding mass \(M\) retains comparable mass and has contributing
prior weight at most \(KM/\rho\).
Only the first deletion removes fractions of the fine states.
Its loss thresholds can be made small enough, in sum over the
finite grid, to retain the spanning probability there. Since a
fine state determines its strip and base \(z_p\)-cell, mixtures
in a base cell also have a fixed positive probability of a
determinant at least \(K^{-1}\).

Apply \Cref{a05:determinant-kakeya} in \(\R^m\) to the full affine
base lines in a strip. If \(m_E\) are the retained masses in its
base cells, the duration in a cell is at most \(Kp\), and hence
\[
 \sum_E\big((m_E/p)^m/K\big)^{1/(m-1)}
       \leq K(M/\rho)^{m/(m-1)}.
\]
At least \(K^{-1}(\rho/p)^m\) base cells are occupied.
Each corresponds to a physical \(W_p\)-cell, and each physical
cell meets at most \(K\) strips, since \(|G_C|\leq K\).
A region in \(C\) with normal coordinate in an interval of
length \(O(Kr)\), \(p\leq r\leq\rho\), is covered by
\(K(\rho/r)^m\) cubes of side \(r\).
Its physical support count is at most
\(K(\rho/r)^m(r/p)^{m+d}\).
Division by the filling per strip proves
\begin{equation}\label{a05:strip-count}
 \#\{\text{occupied }p\text{-strips meeting a normal }r
                   \text{-interval}\}\leq K(r/p)^d.
\end{equation}
Use each filling only to certify its upper support list.
Those lists persist on further restrictions, and intermediate
resolutions follow by subdivision and enlargement.

\paragraph{The parameter zoom.}
The physical cubes \(C\) were used to count scalar-normal strips.
To obtain the conditional spatial mass bounds, we now condition on
a full coefficient bin \(B_\rho\) and rescale its trajectories.
Run the following construction in the family of worlds
\(\mathfrak q=(\gamma,B_\rho)\), where \(\gamma\) is a parent
and \(B_\rho\) is a full parameter block of positive
\(\nu^\circ_\gamma\)-mass. Coordinates, cells and rows in this
construction always include \(\mathfrak q\) in their labels.
Within one such world subtract its central coefficient curves
\(X_c^{\mathfrak q}(t),y_c^{\mathfrak q}(t)\) and divide positions
by \(\rho\):
\[
 x=\frac{X-X_c^{\mathfrak q}}{\rho},\qquad
 \widetilde y=\frac{y-y_c^{\mathfrak q}}{\rho}.
\]
Time is unchanged. All new trajectory coefficients are bounded;
aligned width-\(r'\) parameter bins are precisely the subbins
\(B_{\rho r'}\) of \(B_\rho\).
Choose a time-only row \(F_{\mathfrak q}(t)\) representing
\(G_y\) at active points in each sufficiently fine time bin of
length at most \(\rho\), and extend it boundedly at inactive times.
For the within-world calculations we suppress the superscript on
the central curves and write \(F=F_{\mathfrak q}\).
From \Cref{a05:field-lipschitz},
\begin{equation}\label{a05:zoom-row}
 |F(t)-G_y(W_b)|\leq K\rho,\qquad
 |F(t)-F(u)|\leq K(|t-u|+\rho)
\end{equation}
at active occurrences. The physical positions lie within \(K\rho\)
of a central curve of speed \(K\).
At fixed time, \(\rho(x-F(t)\widetilde y)\), up to a common
translation in each of the at most \(K\) nearby cubes \(C\),
differs from the strip coordinate by at most \(K(\rho^2+b)\).
This is negligible compared with \(\rho\zeta\).
Thus \Cref{a05:strip-count} supplies the scalar-normal local
\(d\)-support bound down to \(\zeta\), also on enlarged upper
intervals by at most \(K\) subdivisions.

Take the fixed prior \(\nu^\circ(\,\cdot\mid B_\rho)\), retaining
all original trajectory tags. At the required resolutions
\(\zeta\leq q'\leq r'\leq1\), for a participating subbin
\(B_{\rho r'}\), the fixed-time numerator in a zoomed spatial
\(q'\)-cube is at most
\[
 K(\rho q')^{j+d}\pi([V]_{\rho r'}).
\]
Its prior denominator is at least
\((\rho r')^{j+d}\pi([V]_{\rho r'})/K\), by
\Cref{a05:parameter-floor}. Consequently the conditional spatial
cap is
\begin{equation}\label{a05:zoom-conditional-cap}
 K(q'/r')^{j+d}.
\end{equation}
The case \(r'=1\) gives the unconditional spatial cap.
These are caps for the original parameter-conditioned prior.
They are not renormalized when labels are subsequently restricted.

\paragraph{The residual velocity.}
Put \(R=x'-F(t)\widetilde y'\).
Within a fixed \(W_b,B_\rho\), its range is much smaller than
\(\zeta\). Indeed, writing \(v_c=y_c'\), one has
\[
 \rho R
 =G_0+G_yv_c-X_c'(t)
       +(G_y-F(t))(V-v_c)+O(b^\kappa).
\]
Here \(|V-v_c|\leq K\rho\), \(|X_c''|\leq K\), and the time
variation within \(W_b\) is at most \(b\).
The choices of \(b,\rho,\zeta\) give the asserted range.

Fix \(s'\geq\zeta\) and put \(\beta=\rho s'\).
A \(\beta\)-parameter block uses at most \(K/s'\) zoomed
\(s'\)-cells in its unique \(\rho\)-zoom.
Delete block/cell entries of incidence mass below
\(s'\nu^\circ(B_\beta)/K_1\).
The summed loss is at most \(K/K_1\), before any normalization
inside \(B_\rho\).
In one zoomed cell the time interval has length \(s'\), whereas
the physical spatial cross-section at each time has side
\(K\beta\) and moves with the central curve.
For a fixed velocity bin \(J=[V]_\beta\), integrating
\Cref{a05:fixed-time-cap} along that moving cross-section gives
\[
 \lambda(\text{zoomed cell},J)
       \leq Ks'\beta^{j+d}\pi(J).
\]
Every surviving block pays at least
\[
 \frac{s'\nu^\circ(B_\beta)}{K_1}
       \geq K^{-1}s'\beta^{j+d}\pi(J).
\]
Both estimates contain \(s'\), so there are at most \(K\)
blocks. No union over physical cells of time width \(\beta\)
is used. Within one block the zoom coefficients have diameter
\(O(s')\); their absolute sizes are \(K\).
Across the cell's time interval, \(x'\) varies by \(Ks'\), and
\Cref{a05:zoom-row} varies by at most \(K(s'+\rho)\leq Ks'\).
Thus one block contributes at most \(K\) width-\(s'\) residual
bins for \(R\).

We next remove the velocity condition with a summed denominator
argument. Write \(K_{\rm pre}=N^{o(1)}\) for the accumulated bounds, and let \(m_0(W)\) be the earlier unnormalized fine-state
weights, including mixture weights, and retain
\[
 \lambda(W,J)\leq K_{\rm pre}m_0(W)\pi(J),\qquad
 \sum_Wm_0(W)\leq K_{\rm pre}.
\]
For a coarse velocity bin \(I=[V]_\rho\), the parameter list
allows \(K_{\rm pre}\) blocks \(B_\rho\) given \(W,I\), and each has
\(K_{\rm pre}\) zoom cells meeting \(W\), since \(K b/\rho\ll\zeta\).
There are at most \(K_{\rm pre}^2\) refined entries
\(E=(W,B_\rho,s'\text{-cell})\) over \(W,I\).
Deleting those of mass below \(m_0(W)\pi(I)/K_{\rm cut}\) loses at most
\[
 \frac{K_{\rm pre}^2}{K_{\rm cut}}\sum_Wm_0(W)\sum_I\pi(I)
 \leq \frac{K_{\rm pre}^3}{K_{\rm cut}}.
\]
Choose the subpower \(K_{\rm cut}\) larger than all accumulated inverse-subpower
density losses. On each survivor,
\begin{equation}\label{a05:zoom-posterior}
 \lambda(J\mid E)\leq K\frac{\pi(J)}{\pi(I)}.
\end{equation}
This uses neither a new prior nor a current incidence floor from
\Cref{a05:parameter-floor}.

Every occurring residual bin has a retained entry witnessing it.
Near constancy within that entry and \Cref{a05:zoom-posterior}
make its supporting velocity bins have total palette mass at
least \(\pi(I)/K\), within boundedly many neighboring residual
bins. A fixed fine velocity bin supports at most \(K\) residual
bins by the preceding block/cell count. Summing against the
common palette in \(I\) gives at most \(K\) residual bins in
the zoomed cell.

\paragraph{Residual comparison and the scalar-normal input.}
It remains to make a common residual choice in each zoomed cell.
Carry out the preceding deletions on a fixed finite grid
\(\{s_1,\ldots,s_L\}\subset[\zeta,1]\), including its endpoints
up to dyadic rounding. Choose the thresholds successively so that
the summed loss is at most half the entering dense mass. The two
losses at a tested width have the forms \(K/K_1\) and
\(K_{\rm pre}^3/K_{\rm cut}\) obtained above, so on a fixed grid
the required thresholds remain subpower. Denote the remaining
unnormalized measure, with the original parent weights, by \(\mu_0\).
Its mass \(\delta_0\) is at least \(K^{-1}\).

In world \(\mathfrak q=(\gamma,B_\rho)\) use
\[
 Z_{\mathfrak q}(e)=(t,\widetilde y_l(t),x_l(t)),\qquad
 R_{\mathfrak q}(e)=x_l'(t)
                  -F_{\mathfrak q}(t)\widetilde y_l'.
\]
The dimension of \(Z_{\mathfrak q}\) is \(D=j+2\), and time retains
its parent units. Freeze the width-\(s_i\) residual lists in every
\(s_i\)-cell of these coordinates. The preceding argument bounds
their sizes by \(K_i=N^{o(1)}\). The palette paying for each list is
common to the cell because \(B_\rho\) fixes \(I=[V]_\rho\).
The finer entry \((W,B_\rho,C)\) supplied a denominator and a witness
for this list; it is not an additional key for the residual choice.

Apply \Cref{a05:common-local-choice} with \(Y=R_{\mathfrak q}\).
Choose the slowly enlarging grid together with the preceding
deletions and density losses as specified after that lemma. We obtain
a restriction \(\mu_L\) of mass \(\delta\geq K^{-1}\) such that
\begin{equation}\label{a05:zoom-residual-comparison}
 |R_{\mathfrak q}(e)-R_{\mathfrak q}(\widehat e)|
 \leq K\bigl(|Z_{\mathfrak q}(e)-Z_{\mathfrak q}(\widehat e)|
                         +\zeta\bigr)
\end{equation}
for every pair of retained incidences in the same world.
Coordinates and rows need not agree across worlds. This time the
selection may remove individual edges: the last use of breadth in
this construction was strip filling, and its remaining steps use only the
inherited upper bounds.

Put \(a_{\mathfrak q}=p_\gamma\nu^\circ_\gamma(B_\rho)\).
These weights sum to at most one. Hence some positive-weight world
satisfies \(\mu_L(\mathfrak q)/a_{\mathfrak q}\geq\delta\).
Divide by this parent weight and parameter-block mass. Since
\(\lambda\leq K\nu^\circ\,dt\), the retained labels have
product-prior mass at least \(K^{-1}\) for the fixed probability
\(\nu^\circ_\gamma(\,\cdot\mid B_\rho)\,dt\).

We can now apply the scalar-normal results of \Cref{sec:scalar-tools}
in this world. The trajectories are \((x,\widetilde y)\), the row is
\(F_{\mathfrak q}\), and the prior remains
\(\nu^\circ_\gamma(\,\cdot\mid B_\rho)\). The strip count gives
the scalar-normal support bound; \Cref{a05:zoom-conditional-cap}
gives both spatial caps, including conditioning on a parameter bin;
\Cref{a05:zoom-row} gives row comparison; and
\Cref{a05:zoom-residual-comparison} gives residual comparison.
These are exactly the hypotheses of \Cref{a04:scalar-data},
with the floor already chosen for \Cref{prop:higher-jets}.
Apply \Cref{prop:higher-jets}. It produces a
degree-\(M\) polynomial with coefficients at most \(K\) fitting \(F\),
hence \(G_y\), to \(Kh^{M+1/4}\) along labeled portions of an indexed
line of duration at least \(h/K\), including the counted occurrences.
Use \Cref{a05:velocity-graph-estimate} and
\(G_0=X'_l-G_yV_l+O(b^\kappa)\) to fit \(G_0\) too.
All witness times on such a line in the block have the same true
\(Q_s\), by its block-invariant trajectory assignment.
Write \(K'\) for the combined subpower in this discovery.
For the fixed \(\epsilon>0\) and all sufficiently large \(N\),
the coefficient bound is smaller than \(N^\epsilon\), the
duration is at least \(hN^{-\epsilon}\), and
\(K'h^{M+1/4}\leq h^M/4\). Coefficient rounding still leaves
error below \(h^M\) and coefficients below \(N^\epsilon\).
These reference-labeled witnesses qualify occurrences lying in
\(U_{N,\epsilon}\), a contradiction. The discovered mass can be
inverse-subpower: any positive mass contradicts the definition
of this failure set. Therefore, for every fixed \(\epsilon>0\),
\[
 \frac{\Lambda_{\rm ref}(U_{N,\epsilon})}{L_{\rm ref}}
       \longrightarrow0.
\]

For \(\epsilon_r=1/r\), choose cutoffs beyond which this fraction
is at most \(1/r\), and let \(r(N)\to\infty\) slowly enough to
respect those cutoffs and the preceding finite-grid preparations.
Set \(K_{\rm qual}=N^{1/r(N)}\) and
\(U_N=U_{N,1/r(N)}\). Then \(K_{\rm qual}\) is subpower and
\Cref{a05:qualification-exception} follows. For any required fixed
finite collection of applications, divide the failure budget
among its members before enlarging the collection. This proves a
relative \(o(1)\) exception in the fixed substantial mixture;
it gives neither finite-\(N\) nullness nor a rate that can be
passed to an arbitrary later sparse restriction.
\end{proof}


\subsection{Cartesian interpolation of the physical field}

Put \(p_*=1/k\), so \(a\sim_{\log,N}h^{p_*}\), and choose
\begin{equation}\label{a05:improvement-scales}
 \begin{gathered}
 0<\alpha<\tfrac12,\qquad
 \eps_T=h^{p_*-1+\alpha},\qquad
 a_1=ah^{e_1},\qquad e_1=\alpha/2,\\
 c_0=e_1(1-d)/2,\qquad
 a_2=ah^{e_2},\qquad e_2=c_0/3.
 \end{gathered}
\end{equation}
In particular \(0<e_2<1/24\).
Take a dyadic mesh \(\xi\) in block-normalized base coordinates so
fine that all later evaluations are accurate, in particular
\(Kh\xi\ll a_1\). Velocity bins are still finer.
The physical mesh \(b\), the polynomial-qualification meshes,
and the graph preparations at depths of thickness \(a_1,a_2\)
were chosen with these fixed power margins before constructing
\(G\). The displayed exponents determine the requirements in
advance; none depends on the magnitude of the hidden coordinate.

Use the states
\[
 \mathcal D=(\text{parent world},Q_s,W_b).
\]
They have actual list knowledge and the required deep-history
ambiguity. Return to the substantial graph-reference measure
\(\Lambda_{\rm ref}\). Prepare conditional palette domination
and breadth here, with
strip probabilities as small a fixed constant as needed for all
subsequent fixed-size tuples and fractional losses.
Denote this substantial submeasure by
\(\Lambda_{\rm prep}\leq\Lambda_{\rm ref}\). For every state of
positive prepared mass, put
\[
 m_{\mathcal D}=\Lambda_{\rm prep}(\mathcal D),\qquad
 q_{\mathcal D}
 =\frac{\Lambda_{\rm prep}(U_N\cap\mathcal D)}{m_{\mathcal D}}.
\]
For any fixed \(\eta>0\), \Cref{a05:qualification-exception} gives
\begin{equation}\label{a05:good-state-deletion}
 \sum_{q_{\mathcal D}>\eta}m_{\mathcal D}
 \leq\eta^{-1}\Lambda_{\rm prep}(U_N)
 \leq(e_N/\eta)L_{\rm ref}=o(1).
\end{equation}
Thus the states with \(q_{\mathcal D}\leq\eta\) retain substantial
mass. If a spanning \(m\)-tuple has probability at least
\(p_0>0\) in each prepared state, choose the fixed threshold
\(\eta\leq p_0/(2m)\). A union bound leaves probability at least
\(p_0-m\eta\geq p_0/2\) for a spanning tuple all of whose
occurrences are qualified. These are the good states used below.
All their prepared reference edges remain available. Qualification
will supply the interpolation witnesses; after proving a statement
about a state, we may restore the other prepared edges of that
state without claiming they too are qualified.

Choose a parent and \(Q=Q_s\) with good-state mass at least
\(h\nu_Q/K\). Such a choice exists since the sum of \(h\nu_Q\)
over packets, blocks and parent mixture weights is at most one.
Write
\[
 u=(z-z_c)/h
\]
with a fixed origin such that \(|u|\leq K\) on the whole block;
there is no velocity shear. Dividing incidence by \(h\nu_Q\)
gives a base density at most \(K\) in \(u\)-space.

Over a specified sufficiently fine base cell, \(G\) to width
\(\eps_T\) has only \(K\) possible bins on the preceding reference
used for fit times. Indeed the \(Q\)-graph lists predict \(X\)
to \(Ka\) with physical gradient at most \(K\).
Since \(\eps_T\gg Ka\), a sufficiently small base cell meets at
most \(K\) physical cells at resolution \(\eps_T\).
Apply the coarse field lists in \Cref{a05:velocity-graph}.

Conditional product sampling at a good state gives tuples of
\(m\) very fine velocity bins with qualifying witnesses and
determinant at least \(K^{-1}\) with probability bounded below.
Palette domination makes their product-palette measure at least
\(K^{-1}\). Averaging fixes one tuple with common witnesses on
states of mass at least \(h\nu_Q/K\).
Let \(A_0\) have columns \({\bf v}_i=(1,v_i)\), the bin centers.
Then \(\|A_0\|+\|A_0^{-1}\|\leq K\).
Use a product grid of step \(\xi\) in \(A_0^{-1}u\), assigning
states by their fine-cell centers.

For each direction \(i\) and grid fiber, consider the polynomial
fits from \Cref{a05:qualification} on witnessing lines in that
velocity bin. Their normalized-time coefficient vectors, to
width \(\eps_T\), have at most \(K\) possible bins.
To prove this, choose one fit and one witnessing line for each
represented bin. Such a line stays within \(K\xi\) of the fiber
throughout the block, because its direction bin is finer than
\(\xi\). At each time its base position is known to \(K\xi\);
the time component of every basis direction equals one.
Sample \(M+1\) independent uniform block times.
For each represented bin, the probability that all these times
lie in its witness set and have pairwise gaps at least \(h/K_1\)
is at least \(K^{-1}\), for sufficiently large subpower \(K_1\).
At any successful tuple, the base-only field lists and polynomial
interpolation permit at most \(K\) coefficient bins, since
\(h^M\ll\eps_T\). Summing this probability proves the claim.

For each state fix one qualifying witness, with its polynomial fit,
in each of the \(m\) selected velocity bins. Independently choose a
polynomial bin on every fiber and one field bin at every grid cell,
uniformly from their respective lists.
Retain a state when its row \(G(W_b)\) lies in the chosen field bin
and all its \(m\) witness fits lie in the chosen polynomial bins on
the corresponding fibers. These \(m+1\) tests involve lists of size
at most \(K\), so some choices retain mass at least \(h\nu_Q/K\).
Use the center of the selected field bin as the chosen cell value.
It is within \(K\eps_T\), on every chosen fiber, of one polynomial
of degree at most \(M\) in that fiber coordinate. Indeed, use
representative coefficients from the selected polynomial bin:
normalized time is affine in that coordinate, and its discrepancy
from witness time is at most \(K\xi\).

Delete cells of mass below \(h\nu_Q\xi^m/K_1\).
There are at most \(K\xi^{-m}\) possible cells, so the loss is
small. The base ceiling shows that the remaining subset occupies
a fraction \(\eta_0\geq K^{-1}\) of the product grid.
Choose \(M+1\) anchor coordinates and one target coordinate
independently on each axis. The probability that all their
product vertices are retained is at least
\(\eta_0^{(M+2)^m}\).
For completeness, condition on the coordinates on all but one
axis. Replacing one coordinate by several independent coordinates
raises its conditional success probability to a fixed power.
Jensen's inequality bounds the average of that power below by
the same power of the preceding average. Iterating over the
axes yields the stated bound.
The probability of any two anchors on an axis lying within
\(K_1^{-1}\) is smaller than this bound when \(K_1\) is a
sufficiently large subpower; the mesh is a fine fixed power.
Fix separated anchors with an inverse-subpower fraction of
successful targets.

Tensor Lagrange interpolation from the pure anchor vertices gives
a polynomial row \(T(u)\), of degree at most \(M\) separately in
each product coordinate and bounded total degree, with coefficients
and norms at most \(K\).
At a successful target, interpolate one coordinate at a time
using the chosen fiber polynomial. All necessary intermediate
vertices are present, and inverse-subpower anchor gaps bound the
Lagrange weights by \(K\). The error is at most \(K\eps_T\).
The per-cell mass floors convert the target count back to mass
at least \(h\nu_Q/K\). Transferring from centers to events costs
negligibly more:
\begin{equation}\label{a05:interpolated-field}
 |G(W_b)-T(u)|\leq K\eps_T.
\end{equation}
This statement depends only on the state. We now restore all
prepared reference edges of the selected good states, preserving
their breadth and palette bounds.

For a line of weighted duration at least \(h/K\) in this family,
\Cref{a05:velocity-graph-estimate,a05:interpolated-field} and
polynomial Remez in normalized block time give
\begin{equation}\label{a05:line-field-comparison}
 |X_l'(t)-T(u_l(t))\cdot{\bf v}_l|\leq K\eps_T
 \quad\text{throughout its }Q\text{-block}.
\end{equation}
The comparison is polynomial in time with a fixed degree.

We will also need the joint support count
\begin{equation}\label{a05:joint-support}
 \#\{(\xi\text{-bin of }u,\ a_i\text{-bin of }X)\}
       \leq K\xi^{-m}(a/a_i)^d,\qquad i=1,2.
\end{equation}
Use true descendants at thickness \(a_i\), with horizons
\(h_i\sim_{\log,N}a_i^k\), rounded with \(K\) loss.
They have trajectory mass at least \(a_i^dh_i^j/K\) and are
disjoint subsets of \(Q\)'s trajectories per \(h_i\)-block.
Each uses at most \(K(h_i/(h\xi))^m\) joint cells by its bounded
gradient graph list; \(\xi\) is chosen sufficiently fine.
There are \(h/h_i\) such time blocks. Thus their total count is
\[
 K\frac{h}{h_i}\frac{\nu_Q}{a_i^dh_i^j}
       \left(\frac{h_i}{h\xi}\right)^m
 =K\xi^{-m}\frac{\nu_Q}{a_i^dh^j}
 \leq K\xi^{-m}(a/a_i)^d.
\]

\subsection{Switched lines force approximate conservativity}

A conservative field has a simple geometric property: its integral along
any broken base path depends only on the two endpoints. If the skew
part of \(DT\) is nonzero, two orders of travel can instead reach the
same base point with different accumulated values of \(X/h\);
\Cref{fig:isotropic-switching} illustrates this local obstruction.
The proof below uses two bases of directions to turn it into a
nonzero Jacobian minor. Its lower bound makes the probability of
landing in the support from \Cref{a05:joint-support} tend to zero
when \(d<1\), contradicting the positive probability of the
switched paths.

\begin{figure}[htbp]
\centering
\begin{tikzpicture}[scale=0.95,>=Stealth,font=\small]
  \coordinate (O) at (0,0);
  \coordinate (A) at (3,0.45);
  \coordinate (B) at (0.8,2.1);
  \coordinate (C) at (3.8,2.55);
  \draw[->,thick,blue!70!black] (O) -- (A)
    node[midway,below] {\(s\mathbf v_i\)};
  \draw[->,thick,blue!70!black] (A) -- (C)
    node[midway,right] {\(t\mathbf v_j\)};
  \draw[->,thick,dashed,orange!80!black] (O) -- (B)
    node[midway,left] {\(t\mathbf v_j\)};
  \draw[->,thick,dashed,orange!80!black] (B) -- (C)
    node[midway,above] {\(s\mathbf v_i\)};
  \fill (O) circle (1.8pt) node[below left] {\(u_0\)};
  \fill (C) circle (1.8pt) node[above right] {\(u_0+s\mathbf v_i+t\mathbf v_j\)};
  \node at (1.9,-0.75) {same base endpoint};
  \draw[->] (7.1,-0.25) -- (7.1,2.85)
    node[above] {accumulated \(X/h\)};
  \fill[blue!70!black] (7.1,2.0) circle (2.2pt);
  \fill[orange!80!black] (7.1,0.85) circle (2.2pt);
  \draw[<->] (7.5,0.85) -- (7.5,2.0)
    node[midway,right] {difference detects \(\operatorname{skew}DT\)};
  \node[align=center] at (7.1,-0.75) {possibly different\\integrals of \(T\cdot du\)};
\end{tikzpicture}
\caption{Two orders of the same two base displacements. The difference
between their line integrals has mixed term
\(st\bigl(DT(u_0)[\mathbf v_i]\cdot\mathbf v_j-DT(u_0)[\mathbf v_j]\cdot\mathbf v_i\bigr)\).
This is a schematic illustration of the obstruction; the proof uses
\(2m\) successive legs, with the first and last \(m\) directions
forming bases.}
\label{fig:isotropic-switching}
\end{figure}

\begin{proposition}[Conservativity of the interpolated field]
\label{a05:conservativity}
In coefficient norm, or equivalently in sup norm on the required
enlarged box up to \(K\),
\[
 \|\operatorname{skew}DT\|\leq K h^{p_*-1+c_0}.
\]
\end{proposition}

\begin{proof}
Otherwise restrict to a subsequence where the skew coefficient
norm is at least \(h^{p_*-1+c_0}\).
Normalize the restored state-selected incidence measure in \(Q\)
to a probability \(\Pp\). Its mass denominator is at least
\(h\nu_Q/K\), so its base density is at most \(K\).
After discarding an arbitrarily small fraction of starting events,
the polynomial sublevel estimate gives skew at least
\(h^{p_*-1+c_0}/K\) at every retained start. A sufficiently large
subpower threshold makes this discarded mass as small as needed.

Starting from \(e_0=(l_0,t_0)\), perform \(2m\) switches.
From \(e_{i-1}\), draw \(e^\diamond=(l_i,t^\diamond)\) afresh
from \(\Pp(\,\cdot\mid\mathcal D(e_{i-1}))\).
Then retain its line and draw \(t_i\) from the conditional
\(\Pp\)-law given \(l_i\), producing \(e_i=(l_i,t_i)\).
Both operations preserve the marginal \(\Pp\).
A line with pre-normalization weighted duration below \(h/K_1\)
has \(\Pp\)-probability at most \(K/K_1\), by the full assigned
prior weight \(\nu_Q\) and the preceding mass denominator.
A union bound therefore removes short choices at arbitrarily
small path cost.

Before these longness restrictions, the velocity at each switch,
conditional on its whole past, is sampled from a prepared state
law. Its probability in any specified thin pencil is as small
as needed. Sequential avoidance therefore gives, with probability
bounded below, a well-conditioned basis from the first \(m\)
velocities and another from the last \(m\).
Take sufficiently fine bins so their centers retain the determinant
bounds. Consequently paths with a high-skew start, these two
bases, and only long choices have probability bounded below.

For every current edge and past history, the \emph{unnormalized}
next-step measure, restricted to long choices, obeys
\begin{equation}\label{a05:transition-domination}
 d\Pp\{[V_{l_i}]=J,\ t_i\in dt_i,\ l_i\text{ long}
                  \mid\text{current edge and past}\}
       \leq K\pi(J)\,\frac{dt_i}{h}.
\end{equation}
The state law supplies \(K\pi(J)\). On each long line,
\Cref{a05:incidence-density} divided by its duration bounds the
endpoint-time density by \(K/h\). We do not normalize the good
transition kernel. Iterating conditional integration backwards
therefore bounds every nonnegative function of the start,
control bins and endpoint times by \(K\) times the product
of the starting \(\Pp\)-law, the palette measures and \(dt_i/h\).
No independence after conditioning on a successful path is needed.
For fixed start, the change to signed increments
\[
 \delta_i=(t_i-t_{i-1})/h
\]
is triangular with Jacobian one in normalized times; they range
in \([-1,1]^{2m}\).

For binned control directions \(\widetilde{\bf v}_i\), simulate
base legs \(\widetilde{\bf v}_i\delta_i\) in \(u\), and simulate
increments in \(X/h\) by integrating
\(T\cdot\widetilde{\bf v}_i\) on these signed legs.
Switching within \(W_b\) introduces only \(Kb\) physical error
when the new line is evaluated at the previous time.
By \Cref{a05:line-field-comparison} and the bounded coefficients
of \(T\), the final simulation errors are
\[
 K(b/h+p_v)\ll\xi
 \quad\text{in }u,\qquad
 K(b/h+p_v+\eps_T)\ll a_1/h
 \quad\text{in }X/h,
\]
where \(p_v\) is the velocity-bin width.
The second comparison has a fixed power margin:
\[
 \eps_T/(a_1/h)\sim_{\log,N}h^{\alpha/2}.
\]

Write the endpoint map as
\[
 \delta\longmapsto
       (u_0+A\delta,\ \mathcal R(\delta)),
\]
where \(A\) has columns \(\widetilde{\bf v}_i\) and
\(\mathcal R\) is a polynomial of fixed degree.
Let \(I,J\) be the first and last \(m\) indices.
The columns in \(I\) determine every base displacement. Consequently
the derivative in a direction \(j'\in J\), minus the corresponding
combination of derivatives in \(I\), leaves the base endpoint fixed.
Its effect on the accumulated height is measured by a maximal
Jacobian minor. For \(j'\in J\), the maximal Jacobian minor using \(I\cup\{j'\}\),
divided up to sign by \(\det A_I\), equals
\[
 \partial_{j'}\mathcal R
 -(\nabla_I\mathcal R)^{\rm T}A_I^{-1}A_{j'}.
\]
Differentiate this identity at zero in each coordinate of \(I\).
If every minor had polynomial norm at most \(\eta\), the cross
Hessian would differ by at most \(K\eta\) from
\[
 A_I^{\rm T}H_{\rm sym}A_J,\qquad
 H_{\rm sym}
   =A_I^{-{\rm T}}(D^2\mathcal R(0))_{I,I}A_I^{-1},
\]
a symmetric matrix transported between the two bases.
Here fixed-degree polynomial derivative bounds control a
derivative at zero by the minor norm.
On the other hand, directly differentiating the successive
integrals gives, for \(i\in I,j'\in J\),
\[
 (D^2\mathcal R(0))_{i,j'}
       =DT(u_0)[\widetilde{\bf v}_i]\cdot
                    \widetilde{\bf v}_{j'}.
\]
Later legs have zero duration at zero, so there are no additional
cross terms. Inverting both bases bounds the skew of \(DT(u_0)\)
by \(K\eta\). Hence for the high-skew starts at least one maximal
minor has polynomial norm at least \(h^{p_*-1+c_0}/K\).

For every start and control tuple, choose such a minor and remove
the region where its absolute value is below
\(h^{p_*-1+c_0}/K_1\). Fixed-degree polynomial sublevels in
Lebesgue increments, together with
\Cref{a05:transition-domination}, make the removed path mass
arbitrarily small. On the passing region, fix the complementary
\(m-1\) increment coordinates. The regular fibers of the remaining
polynomial map have bounded multiplicity, by the fixed-degree
semialgebraic fiber bound. The change-of-variables formula thus
bounds the Lebesgue measure mapped into one endpoint cell by
\[
 K\xi^m(a_1/h)\,h^{-(p_*-1+c_0)}.
\]
The complementary coordinates have bounded volume. The resulting
cell estimate has the full \(m+1\)-dimensional volume factor
\(\xi^m(a_1/h)\). It remains to sum it over the smaller support
available to the actual endpoints.
Every actual long path has its simulated endpoint in a bounded
enlargement of the union counted by \Cref{a05:joint-support}.
Summing and integrating the start and palette weights bounds
the successful-path probability by
\[
 \begin{aligned}
 K(a/a_1)^d(a_1/h)h^{-(p_*-1+c_0)}
 &\leq K h^{e_1(1-d)-c_0}\\
 &=K h^{e_1(1-d)/2}\longrightarrow0.
 \end{aligned}
\]
This contradicts the positive lower probability and proves the
proposition.
\end{proof}

\subsection{A potential and concentrated intercept groups}

Define the polynomial
\[
 P(u)=h\int_0^1T(tu)\cdot u\,dt.
\]
Differentiating under the integral and comparing with \(hT(u)\)
leaves only the skew derivative integrated against a bounded
multiple of \(u\). Thus \Cref{a05:conservativity} gives
\begin{equation}\label{a05:potential-error}
 |\nabla P-hT|\leq Kah^{c_0}.
\end{equation}
The norms of all fixed positive derivatives of \(P\) are at
most \(Kh\) on the needed enlarged box.
For a long line in \Cref{a05:line-field-comparison},
the variation of \(X-P(u)\) across the block is at most
\(Kah^{c_0}\): integrate the directional error over time \(h\),
using \(h\eps_T\leq Kah^\alpha\) and \(c_0<\alpha\).

Bin the midpoint value of \(X-P(u)\) to width \(a_2\), denoting
its bin and center by \(C\). This is invariant along the
assigned trajectory. At a state \(\mathcal D\) of long labels,
only boundedly many neighboring \(C\)'s occur:
the variation \(Kah^{c_0}\) is much smaller than \(a_2\), the
state knows \(X,u\) much more accurately, and \(P\) has controlled
gradient.
Within one intercept group, the normalized base density, after
division by \(h\nu_Q\), is at most
\begin{equation}\label{a05:intercept-base-cap}
 K(a_2/a)^d.
\end{equation}
Indeed apply the exact-base \(X\)-cap at width \(a_2\), whose
center follows \(P(u)+C\), then use
\(\nu_Q\sim_{\log,N}a^dh^j\).

We next bound the number of groups without asserting that an
arbitrary restriction preserves breadth. First remove short
lines globally in \(Q\), and states with excessive fractional
loss. In each remaining state remove the fractions belonging
to \(C\)'s of mass below a sufficiently small fixed fraction.
There are only boundedly many \(C\)'s there, so the lost fraction
is small. Remove short lines again in this family, still with
only one \(C\) per line. Remove \((\mathcal D,C)\)-states with
excessive fractional loss, and groups with excessive total
loss compared with before the last short-line deletion.
Mass at least \(h\nu_Q/K\) remains.
If a group has remaining mass \(M\), its contributing prior
weight is at most \(KM/h\), using the immediately preceding
duration test.
At each surviving \((\mathcal D,C)\), the law is at least a
fixed fraction of the original good state law, up to the chosen
small losses. The pencil probabilities were chosen smaller
than these fractions. All needed breadth and palette bounds
therefore survive.

Apply \Cref{a05:determinant-kakeya} in dimension \(m\) to the
affine base lines of one group, in \(u\)-coordinates.
Assign \(\xi\)-cells by state centers. A line spends at most
\(Kh\xi\) weighted physical time in one such cell, and its
radius-\(O(\xi)\) tube contains the assigned cell.
If \(m_E\) are the group cell masses, determinant success and
the prior-weight bound imply
\[
 \sum_E\big((m_E/(h\xi))^m/K\big)^{1/(m-1)}
       \leq K(M/h)^{m/(m-1)}.
\]
Power mean gives at least \(K^{-1}\xi^{-m}\) occupied base cells
per group. Each joint cell at \(X\)-width \(a_2\) pays for at
most \(K\) intercept groups, since the variation of \(P\) on
a base cell is negligible compared with \(a_2\).
The count in \Cref{a05:joint-support} with \(i=2\) now yields
at most \(K(a/a_2)^d\) groups.

Delete groups below
\begin{equation}\label{a05:heavy-group}
 K^{-1}h\nu_Q(a_2/a)^d
\end{equation}
with a sufficiently large denominator. The summed loss is small.
On every retained long line in any group,
\begin{equation}\label{a05:potential-fit}
 |X_l(t)-P(u_l(t))-C|\leq Ka_2
 \quad\text{throughout the entire }Q\text{-block}.
\end{equation}
All subsequent counts use these heavy groups and their immediately
preceding long witnesses.


\subsection{Seven directions and the exceptional conic}

The remaining task is to replace the potential fit by an allowed
quadratic test. In three base variables, even directions with every
triple independent can all lie on a nonsingular conic. Third
directional tests along those directions therefore need not force
every higher-order term to vanish. We must identify the cubic and
quartic terms that survive, and later express them using the hidden
coordinate so that the final test is quadratic.

For a symmetric rank-\(r\) tensor \(T\) on \(\R^3\), associate the
homogeneous form \(F_T(x)=T(x,\ldots,x)\).
Testing \(T(p_i,p_i,p_i,\cdot,\ldots,\cdot)=0\) is equivalent to
the vanishing at \(p_i\) of all partial derivatives of \(F_T\)
of order \(r-3\), up to nonzero constants depending only on \(r\).
We call these the three-copy tests. For seven directions with
every triple independent, the following lemma identifies their
kernels. If the directions lie on a conic \(q=0\), the rank-three
and rank-four kernels are \(q\) times a linear form and a multiple
of \(q^2\), while the rank-five kernel vanishes.
The cubic two-copy test \(T(p_i,p_i,\cdot)\) has a
different role: it is injective even in the conic case and will
control the third derivatives of a branch from its differentiated
second-directional tests. The following lemmas also supply the
conditioning bounds and the passage from approximate conic
directions to exact real ones needed for those applications.

\begin{lemma}[Seven-direction tensor tests]\label{lem:tensor-kernels}
Let \(p_i=(1,v_i)\in\R^3\), \(1\leq i\leq7\), with every triple
linearly independent.
\begin{enumerate}
\item The space of homogeneous quadratics vanishing at all \(p_i\)
has dimension at most one; every nonzero member is nonsingular.
\item The map on symmetric cubic tensors
\[
 T\longmapsto\bigl(T(p_i,p_i,\cdot)\bigr)_{i=1}^7
\]
is injective.
\item If a nonzero quadratic \(q\) vanishes at all seven points,
the kernels of the three-copy tests at ranks \(3,4,5\) are,
under the homogeneous-form correspondence,
\[
 q\,(\R^3)^*,\qquad \R q^2,\qquad \{0\},
\]
respectively.
\end{enumerate}
These assertions have polynomial conditioning. More precisely, if
\[
 |p_i|\leq A,\qquad
 |\det(p_i,p_j,p_k)|\geq A^{-1},\qquad A\geq2,
\]
the inverse of the injective map in part~2 costs at most \(CA^C\).
On the exact-conic family, with the norm, inverse norm and inverse
determinant of \(q\) also bounded by fixed powers of \(A\),
distance to each kernel in part~3 is at most \(CA^C\) times
the norm of the corresponding tests. The constants depend only
on these fixed finite-dimensional spaces and on the fixed
conditioning powers.
\end{lemma}

\begin{proof}
All divisibility arguments can be made over \(\C\).
A complex line contains at most two of the seven points, since
three real linearly independent vectors are also complex linearly
independent. A singular projective conic, being a union of two
lines with multiplicity allowed, cannot contain five of them.

There is exactly one conic through any five, up to scalar.
Indeed, use the first three points as a projective coordinate
basis. Quadratics vanishing there have form
\(a x_0x_1+b x_0x_2+c x_1x_2\).
The remaining two points have all coordinates nonzero.
Their vectors \((x_0x_1,x_0x_2,x_1x_2)\) are not proportional,
since their ratios recover the projective point.
They give two independent conditions on \((a,b,c)\).
Thus the five-point conic space is one-dimensional, and its
nonzero member is nonsingular by the preceding paragraph.
Part~1 follows.

If \(T(p_i,p_i,\cdot)=0\) for a cubic form \(F\), all its
first partial derivatives belong to that quadratic space.
If the space is zero, \(F=0\). Otherwise write
\(\partial_jF=c_jq\). Euler's identity yields \(F=qL\)
for a linear form \(L\). At every tested point,
\[
 0=\nabla F(p_i)=L(p_i)\nabla q(p_i).
\]
Nonsingularity gives \(L(p_i)=0\), and three independent points
force \(L=0\). This proves part~2.

A cubic vanishing at seven distinct points of a nonsingular
conic is divisible by its equation. Here is an elementary
justification of this intersection fact. Over \(\C\), change
coordinates to write the conic as \(x_0x_2-x_1^2=0\), with
parametrization
\([s:t]\mapsto[s^2:st:t^2]\).
The restriction of a cubic is a binary sextic; seven distinct
projective zeros make it identically zero.
For any homogeneous degree, reduction modulo \(x_0x_2-x_1^2\)
leaves monomials using no pair \(x_0,x_2\) simultaneously.
Their substitutions are distinct binary monomials, with the
common monomial \(x_1^r\) counted once. Thus a zero restriction
has zero remainder, proving divisibility.
The rank-three kernel is therefore exactly \(q\) times the
linear forms.

For a quartic \(F\) in the rank-four kernel, each first partial
derivative is such a cubic and is divisible by \(q\).
Euler gives \(F=qH\), with \(H\) quadratic. Then
\[
 \partial_jF=(\partial_jq)H+q\partial_jH.
\]
Choose a nonzero linear partial derivative of \(q\).
The irreducible quadratic \(q\) is relatively prime to it,
so divisibility of \(\partial_jF\) by \(q\) forces \(q\mid H\).
Hence \(F\) is a multiple of \(q^2\).
Conversely, \(q^2\) and its first derivatives vanish at the
seven points.

For a quintic in the rank-five kernel, every second derivative
is a cubic divisible by \(q\). Each first derivative consequently
satisfies the preceding quartic test and is a multiple of \(q^2\).
Euler gives \(F=q^2L\) with \(L\) linear.
Differentiation and reduction modulo \(q^2\) force
\(q\mid(\partial_jq)L\), and hence \(q\mid L\) for a nonzero
partial derivative. Therefore \(L=0\).
This proves all kernel equalities.

For fixed \(A\), the point configurations form a compact
semialgebraic set. The matrix of the cubic two-copy test has
polynomial entries and is injective everywhere on this set,
so its least singular value has a positive minimum.
For the three-copy tests restrict to the compact exact-conic
family, imposing the stated coefficient and determinant bounds.
Their kernel dimensions are constantly \(3,1,0\), so their
least positive singular values also have positive minima.
The minima are semialgebraic functions of \(A\), by real
quantifier elimination; orthogonality to a kernel and unit norm
are fixed-format polynomial conditions.

A positive semialgebraic function of one real parameter has a
reciprocal polynomial lower bound at infinity.
To see the rate relevant here, decompose its graph into finitely
many algebraic branches for large \(A\). On each, a nonzero
relation \(P(A,s)=0\) holds. Remove powers of \(s\) from \(P\).
Its constant coefficient in \(s\) is then a nonzero polynomial
in \(A\), bounded below by a power of \(A\) for large \(A\);
all other coefficients are bounded above by powers of \(A\).
For \(0<s\leq1\), the relation forces \(s\geq cA^{-C}\).
Absorb bounded parameter intervals into the constant.
The claimed inverse estimates follow. Crucially the rank-four
estimate is applied on the exact-conic family: its rank can
increase when the points move off that family.
\end{proof}

\begin{lemma}[Restoration of an approximate real conic]
\label{a05:real-conic}
Suppose the points satisfy the conditioning bounds of
\Cref{lem:tensor-kernels}, and a real quadratic \(q(x)=x^{\rm T}
\mathsf Qx\) has
\(A^{-1}\leq\|q\|\leq A\) and \(|q(p_i)|\leq\epsilon\).
If \(\epsilon\) is smaller than a sufficiently large fixed
reciprocal power of \(A\), then
\[
 |\det\mathsf Q|\geq cA^{-C},
\]
and there are exact real zeros \(p'_i=(1,v'_i)\) of \(q\) with
\[
 |p'_i-p_i|\leq CA^C\epsilon.
\]
For sufficiently small \(\epsilon\), every triple determinant
of the shifted points has absolute value at least \((2A)^{-1}\).
\end{lemma}

\begin{proof}
For coefficient-normalized quadratics the function
\[
 |\det\mathsf Q|+\sum_{i=1}^7|q(p_i)|
\]
has a positive minimum on the compact point family:
a zero would be a nonzero singular quadratic through all seven
points, contrary to \Cref{lem:tensor-kernels}.
The semialgebraic argument in that lemma gives a lower bound
\(cA^{-C}\). Normalize the given \(q\), use the smallness of
its evaluations, then rescale to obtain the determinant bound.

It bounds the smallest singular value of \(\mathsf Q\) below
by \(cA^{-C}\), and hence the homogeneous gradient at each
\(p_i\) below by such a bound. Write
\(g_i=\nabla_vq(1,v_i)\).
Euler's identity reads
\[
 \partial_0q(1,v_i)+v_i\cdot g_i=2q(1,v_i).
\]
Since \(|v_i|\leq A\) and \(q(1,v_i)\) is sufficiently small,
the full-gradient bound implies \(|g_i|\geq cA^{-C}\).
With \(e_i=(0,g_i/|g_i|)\),
\[
 q(p_i+te_i)=q(p_i)+|g_i|t+b_it^2,\qquad |b_i|\leq CA.
\]
If \(q(p_i)\ne0\), take \(t=-2q(p_i)/|g_i|\).
For sufficiently small \(\epsilon\), the quadratic term has
absolute value below \(|q(p_i)|/2\), so this value has the
opposite sign from \(q(p_i)\).
The intermediate value theorem gives a real zero at distance
at most \(2|q(p_i)|/|g_i|\). If \(q(p_i)=0\), no shift is needed.
The first coordinate remains one. Multilinearity shows that
triple determinants change by at most
\(CA^2\max_i|p_i'-p_i|\), proving the final assertion.
\end{proof}

\subsection{Restoring quadratic tests}

\begin{proposition}[Degree restoration]\label{prop:degree-restoration}
The potential fits and heavy intercept groups constructed above
yield a candidate with all the properties in
\Cref{thm:isotropic-improvement}, without reducing the
positive exponent \(e_2\) in \(a_2/a=h^{e_2}\).
\end{proposition}

\begin{proof}
We treat the two versions separately.
\paragraph{Version~1.}
Choose a heavy intercept group. In normalized block time,
\Cref{a05:potential-fit} compares \(P\) with the quadratic
\(X-C\) along each of its lines. Polynomial Remez gives
\[
 |D_{\bf v}^3P|\leq Ka_2
\]
there. Conditional product sampling and palette averaging fix
four separated very fine velocity bins with witnesses on common
states carrying an inverse-subpower fraction of the group.
Use their centers as directions. The locations in a state
differ by at most \(Kb/h\) in \(u\), and fixed derivatives of
\(P\) cost \(Kh\). Finer velocity bins and \(b/h\ll a_2/h\)
make the transferred errors negligible.
By \Cref{a05:heavy-group,a05:intercept-base-cap}, these common
tests hold on inverse-subpower base volume.
Polynomial Remez extends them to the enclosing box.
A binary cubic is determined by its values in four separated
directions \((1,V_i)\); inversion of this Vandermonde system
costs \(K\). Hence \(\|\nabla^3P\|\leq Ka_2\).
The quadratic Taylor polynomial of \(P+C\) fits to \(Ka_2\)
on the block box. The test \(w\) minus this polynomial has hidden
derivative one and zero hidden curvature, and its heavy-group
mass is the required one.

\paragraph{Branches in version~2.}
Here \(w\) may be uncontrolled in tangent units, so we use the
native test rather than assume a bound for \(w\).
In each heavy group consider
\begin{equation}\label{a05:branch-equation}
 P_*(z,w)=P(u)+C,\qquad z=z_c+hu.
\end{equation}
At path observations the residual is \(O(Ka_2)\), the absolute
hidden derivative lies between \(K^{-1}\) and \(K\), and the
hidden curvature is at most \(K\).
The discriminant expression
\[
 \Delta(u)
 =(P_{*,w})^2
       -2P_{*,ww}\bigl(P_*(z,w)-P(u)-C\bigr)
\]
is independent of \(w\), because the equation is quadratic in
\(w\) with constant quadratic coefficient.
It lies between \(K^{-1}\) and \(K\) on the observations.
The heavy-group floor divided by
\Cref{a05:intercept-base-cap} gives inverse-subpower base
volume. Polynomial Remez bounds the coefficients and derivatives
of \(\Delta\) on an enlarged complex box by \(K\).
Thus, about every used point, \(\Delta\) is zero-free on a
complex ball of inverse-subpower radius with a fixed margin.

Each observed \(w_l\) is within \(Ka_2\) of a real simple root
of \Cref{a05:branch-equation}. For a quadratic, choose the root
with the same derivative sign as at \(w_l\); the derivative on
the intervening segment has size at least \(K^{-1}\), after
enlarging \(K\), by the positive discriminant and small residual.
For a linear equation, divide by its nonzero coefficient.
The roots extend holomorphically to the stable balls, by a
holomorphic square root of \(\Delta\), or by linear division.

Choose a lattice of real comparison cores, each inside a
fixed several-fold complex enlargement, with radius common
up to \(K\) among the groups. Assign a core by the
\((\mathcal D,C)\)-center. Since \(Kb/h\) is negligible, it
contains the state's used points with an interior margin.
Boundary points can be assigned in a bounded overlapping cover.
Add the sign of the nearby root.
Discard sign fractions smaller than a sufficiently small fixed
fraction; the core choice itself removes no fraction at a state.
Delete line/intercept/core/sign portions with weighted duration
below \(h/K_1\). Each line has one intercept and at most \(K\)
core/sign choices, so the total loss is at most
\(Kh\nu_Q/K_1\). Remove states with excessive fractional loss.
The fixed-fraction comparisons to the earlier good state laws,
with sufficiently small initial pencil probabilities, retain
general position and palette domination.

There are at most \(K(a/a_2)^d\) group/core/sign choices.
One choice therefore has mass at least
\begin{equation}\label{a05:branch-mass}
 K^{-1}h\nu_Q(a_2/a)^d.
\end{equation}
Let \(U\) be its real inner ball and \(f\) its branch.
Every contributing line has an immediately preceding witness
set, in this same choice, of ordinary length at least \(h/K\),
including its currently counted incidences.
All witnesses lie in \(U\), with \(f\) holomorphic on a
sufficient complex enlargement. The convex cores and that
margin ensure that the line chords used below stay in the
controlled domain.

\paragraph{Analytic estimates before tensor inversion.}
Set
\begin{equation}\label{a05:normalized-test}
 R(u)=P_s(z_c+hu,f(u))/a.
\end{equation}
We will apply the two-copy cubic test to \(\nabla^3f\) and the
three-copy tests to the higher derivatives of \(R\). First we
need analytic bounds on \(R\) that hold before moving any
directions to an exact conic.
At the earlier witnesses, the true packet test obeys
\[
 |P_s(z,w_l)|\leq Ka,\qquad
 |P_{s,w}(z,w_l)|\leq K,\qquad
 |P_{s,ww}|\leq K/a,
 \qquad |w_l-f(u)|\leq Ka_2.
\]
Its exact quadratic Taylor formula in \(w\) gives
\[
 |R(u)|\leq K\bigl(1+a_2/a+(a_2/a)^2\bigr)\leq K.
\]
The base density of this choice is bounded by the original
group ceiling \(Kh\nu_Q(a_2/a)^d\).
Together with \Cref{a05:branch-mass}, this gives real projection
volume at least \(K^{-1}\).
The function \(R\) is holomorphic algebraic of bounded relation
degree. \Cref{lem:algebraic-norm} bounds its norm and every
required fixed derivative by \(K\) on the used interiors,
before any tensor inversion.

Along a contributing line, compare \(f\) with the affine hidden
trajectory \(w_l\), and compare \(R\) with the quadratic
\(P_s(z,w_l)/a\), on its earlier long witness set.
Normalized time on its chord has length at least \(K^{-1}\);
the witness set has inverse-subpower relative length and the
chord has the stated complex margin.
The one-variable algebraic norm bound yields, at the incidences,
\begin{equation}\label{a05:directional-branch}
 |D_{\bf v}^2f|\leq Ka_2,\qquad
 |D_{\bf v}^3R|\leq Ka_2/a.
\end{equation}

Let
\[
 S_0=\inf_{L\ {\rm affine}}\|f-L\|_{C^0(U)}.
\]
For \(S_0>0\), apply \Cref{lem:algebraic-norm} to
\(f-L\) with \(\|f-L\|_{C^0(U)}\leq2S_0\).
All fixed derivatives of order at least two are bounded by
\(KS_0\) on a slightly larger interior. If \(S_0=0\), \(f\)
is affine and these derivatives vanish.

Sample seven velocities in each retained state, sequentially
avoiding all preceding points and joining lines. The prepared
breadth gives a positive probability that all triple determinants
of \({\bf v}_i=(1,v_i)\) have size at least \(K^{-1}\).
Palette domination turns this into a product-palette mass at
least \(K^{-1}\) of eligible bin tuples. Averaging fixes one
tuple with common directional witnesses on states carrying an
inverse-subpower fraction of the current mass.
Using fine bin centers transfers their tests to actual base
points in those states. For \(R\), its preceding derivative
bounds make all mesh errors negligible relative to \(a_2/a\).
For \(f\), the Hessian and higher derivative errors are bounded
by \(KS_0\) times the normalized position and velocity errors.
Choose those errors at most \(h^2/K\), independently of \(S_0\).
Common-state mass and the base ceiling again give base volume
at least \(K^{-1}\). Apply \Cref{lem:algebraic-norm} to the
directional derivative functions themselves; differentiation
of the simple root preserves bounded algebraic relation degree
on the stable domain.
We obtain throughout \(U\), also after the fixed further
differentiations needed below,
\begin{equation}\label{a05:seven-tests}
 |D_{{\bf v}_i}^2f|\leq K(a_2+S_0h),\qquad
 |D_{{\bf v}_i}^3R|\leq Ka_2/a.
\end{equation}
The weaker \(S_0h\) allowance absorbs the \(S_0h^2\) transfer
error and subpower losses.
All ball radii and margins here cost only \(K\).

Differentiate the first tests once and apply the injective
cubic two-copy test in \Cref{lem:tensor-kernels}. Its polynomial
conditioning is subpower, so
\begin{equation}\label{a05:third-f}
 \|\nabla^3f\|\leq K(a_2+S_0h).
\end{equation}
If \(S_0\leq a_2/h\), a quadratic Taylor polynomial \(f_2\)
approximates \(f\) on \(U\) within \(Ka_2\).
The test \(w-f_2((z-z_c)/h)\) is quadratic in parent variables,
has derivative one and curvature zero, and fits the long
witnesses. This includes \(S_0=0\).

\paragraph{The conic case.}
Suppose \(S_0>a_2/h\).
Let \(u_c\) be the ball center and \(L'\) the affine Taylor
polynomial of \(f\) there. Define
\[
 q(x)=\frac{1}{2S_0}\nabla^2f(u_c)[x,x].
\]
Taylor's theorem and \Cref{a05:third-f} give
\begin{equation}\label{a05:f-conic}
 \frac{f(u)-L'(u)}{S_0}
       =q(u-u_c)+O(Kh).
\end{equation}
The relative derivative bound gives \(\|q\|\leq K\).
Conversely the definition of \(S_0\) yields
\[
 1\leq \|q\|_{C^0(U-u_c)}+Kh\leq K\|q\|+Kh,
\]
so \(\|q\|\geq K^{-1}\).
The Hessian tests imply \(|q({\bf v}_i)|\leq Kh\).
Choose one parameter \(A=N^{o(1)}\) controlling all coefficient
norms, direction sizes, and inverse triple determinants.
Since \(h\) is a fixed negative power, \Cref{a05:real-conic}
applies. It gives a determinant lower bound \(K^{-1}\) for
\(q\), and exact real zeros \({\bf v}'_i=(1,v_i')\) with
\[
 |{\bf v}'_i-{\bf v}_i|\leq Kh,
 \qquad
 |\det({\bf v}'_i,{\bf v}'_j,{\bf v}'_k)|\geq K^{-1}.
\]
Thus even a nearly singular or nearly definite proposed
quadratic is quantitatively controlled before inversion.

The independent derivative bounds for \(R\) from
\Cref{a05:normalized-test} show that these shifts alter its
rank-three, rank-four and rank-five contractions by at most
\(Kh\). This is negligible compared with
\(a_2/a=h^{e_2}\), because \(e_2<1\).
The differentiated tests in \Cref{a05:seven-tests} therefore
remain \(Ka_2/a\) in the shifted directions.
Apply the exact-conic kernels of \Cref{lem:tensor-kernels}.
The fifth derivatives of \(R\) are bounded by \(Ka_2/a\);
at \(u_c\), the cubic and quartic Taylor tensors are within
\(Ka_2/a\) of \(q\) times a linear form and a multiple of
\(q^2\), respectively. Taylor expansion through degree four gives
\begin{equation}\label{a05:R-restoration}
 R(u)=R_2(u)+q(u-u_c)L_1(u)
                +\lambda_0q(u-u_c)^2+O(Ka_2/a),
\end{equation}
where \(R_2\) is quadratic, \(L_1\) is affine, and all displayed
coefficients are bounded by \(K\).
To justify division in these factors, the coefficient norm of
a product of polynomials of fixed degrees is bounded below by
a positive constant times the product of their coefficient
norms. Normalize the factors and use compactness and the
absence of zero divisors. Since \(\|q\|\geq K^{-1}\), division
by \(q\) or \(q^2\) costs only \(K\).
A constant term of \(L_1\) can be absorbed into \(R_2\).

On the earlier long witnesses, put
\[
 H'(u,w)=\frac{w-L'(u)}{S_0}.
\]
Because \(a_2/S_0<h\), \Cref{a05:f-conic} and
\(|w-f(u)|\leq Ka_2\) imply
\(H'=q(u-u_c)+O(Kh)\) and \(|H'|\leq K\).
Define
\begin{equation}\label{a05:new-test}
 P_{\rm new}(z,w)
 =P_s(z,w)
   -a\left(R_2(u)+H'(u,w)L_1(u)
                         +\lambda_0H'(u,w)^2\right),
 \qquad u=(z-z_c)/h.
\end{equation}
Every term has total degree at most two in \((z,w)\):
\(u,H',L'\), and \(L_1\) are affine in the indicated variables.
Substituting \(H'\) for \(q\) costs \(Kah\), which is below
\(a_2\) by \(h^{1-e_2}\).
Replacing \(f\) by the actual \(w\) in \(P_s\) costs at most
\[
 K|w-f|+(K/a)|w-f|^2
       \leq Ka_2+Ka_2^2/a\leq Ka_2.
\]
Together with \Cref{a05:R-restoration}, this proves the
\(Ka_2\) value bound on the long witnesses.

The exact derivative perturbations are
\[
 \partial_w\bigl(P_s-P_{\rm new}\bigr)
       =\frac a{S_0}(L_1+2\lambda_0H'),\qquad
 \partial_{ww}\bigl(P_s-P_{\rm new}\bigr)
       =\frac{2a\lambda_0}{S_0^2}.
\]
Consequently
\begin{align}
 \bigl|\partial_w(P_s-P_{\rm new})\bigr|
   &\leq K a/S_0
     <K h^{1-e_2}\ll K^{-1},\label{a05:derivative-margin}\\
 |\partial_{ww}P_{\rm new}|
   &\leq K/a+K a/S_0^2
     \leq K/a_2+(K/a_2)h^{2-e_2}
     \leq K/a_2.\label{a05:curvature-margin}
\end{align}
The old derivative lower bound at the counted observations
survives \Cref{a05:derivative-margin}; the upper bound holds
on the earlier witnesses. No magnitude bound for \(w\),
\(L'\), or \(S_0\) was used, only the normalized value \(H'\)
and the displayed lower bound on \(S_0\).

\paragraph{The same time block.}
In both cases the test is an allowed quadratic.
On each contributing version~2 trajectory, \(z(t)\) and \(w(t)\)
are affine, so its value is a polynomial of degree at most two
in time and its hidden derivative has degree at most one.
The earlier witness set has physical length at least \(h/K\)
inside the same \(h\)-length \(Q_s\)-block.
After normalizing that block to unit length, polynomial Remez
costs only \(K\). Thus the value and derivative upper bounds
hold throughout this block. The lower derivative is needed
only on the counted incidences and was proved there.
Curvature is constant in \(w\) for these quadratic tests.
The version~1 comparison has the same fixed-degree time property.
Extrapolation is never made from an \(h/K\)-set to a unit
\emph{parent} time interval.

The retained choice has mass
\(h\nu_Q(a_2/a)^d/K\), by \Cref{a05:branch-mass} or the heavy
intercept floor in version~1. Degree restoration makes no
further discard. Dividing by its time length \(h\) gives
\(\nu_Q(a_2/a)^d/K\), exactly the candidate mass requirement.
\end{proof}

\begin{proof}[Conclusion of \Cref{thm:isotropic-improvement}]
Start with any substantial residual of the reference path.
The physical estimates and \Cref{a05:velocity-graph} retain a
substantial reference family. \Cref{a05:qualification} is an
assertion of relative \(o(1)\) failure under its frozen
post-renewal incidence law. It follows by finding qualifying
occurrences in every hypothetical substantial fixed-tolerance
failure set and then taking the slow diagonal.
The fixed-threshold deletion in \Cref{a05:good-state-deletion}
therefore leaves substantial good-state mass. Cartesian
interpolation selects a packet and states of mass
at least \(h\nu_Q/K\), then restores their prepared reference
edges. The conservativity estimate yields the potential and
heavy intercept groups constructed above; \Cref{prop:degree-restoration}
provides the required quadratic candidate. By
\Cref{a05:improvement-scales}, its gain has positive limit
\(e_2ks\). All meshes and degree bounds used to discover it
depend on fixed exponents, while the actual output construction
and its original-scale guards are those of
\Cref{prop:candidate-upgrade}. The criterion therefore gives
the isotropic contradiction on every substantial residual.
\end{proof}

\begin{corollary}[Zero horizon in the scalar model]
\label{a05:scalar-zero-horizon}
For every fixed admissible choice of angular parameters and every
\(d<1\), a version~1 exact problem in \(C(d)\) has
\(H(E)=0\). In particular there is no bad problem in version~1.
\end{corollary}

\begin{proof}
If a bad version~1 problem existed, the critical-path construction
of \Cref{prop:critical-path} would give a tangent problem with
positive \(k\) and precisely the version~1 data above.
\Cref{thm:isotropic-improvement,prop:candidate-upgrade} would
contradict its maximizing-path exclusion.
\end{proof}

This scalar conclusion is established before the finite-narrowness
analysis. It supplies the version~1 input to
\Cref{thm:scalar-mesh-projection}, whose projection labels are then
used in \Cref{sec:finite-narrowness,sec:critical-rates,sec:unbounded-narrowness}.
The unbounded-narrowness case retains its separate limiting construction.
